\documentclass[12pt,a4paper]{article}

\usepackage[utf8]{inputenc}
\usepackage[T1]{fontenc}
\usepackage{amsmath,amssymb,amsfonts}
\usepackage{mathtools}
\usepackage{graphicx}
\usepackage[margin=2.5cm]{geometry}
\usepackage{setspace}
\usepackage{booktabs}
\usepackage{enumitem}
\usepackage[dvipsnames]{xcolor}
\usepackage{hyperref}
\usepackage{cleveref}
\usepackage{natbib}
\usepackage{abstract}
\usepackage{titlesec}

\hypersetup{
    colorlinks=true,
    linkcolor=blue!70!black,
    citecolor=blue!70!black,
    urlcolor=blue!70!black,
    pdfauthor={Hasok Chang},
    pdftitle={Antimines and the Simulated Architecture Confound},
}

\onehalfspacing
\titleformat{\section}{\large\bfseries}{\thesection.}{0.5em}{}
\titleformat{\subsection}{\normalsize\bfseries}{\thesubsection}{0.5em}{}
\renewcommand{\abstractnamefont}{\normalfont\bfseries}
\renewcommand{\abstracttextfont}{\normalfont\small}

\begin{document}

\begin{center}
    {\LARGE\bfseries Antimines and the Simulated Architecture Confound:\\[4pt]
    Why Prompting Cannot Rewrite Physical Law\\[4pt]}

    \vspace{1.2em}

    {\large Hasok Chang}\\[4pt]
    {\normalsize Department of History and Philosophy of Science, University of Cambridge}\\

    \vspace{1em}
    {\small March 2026}
\end{center}
\vspace{0.5em}

\begin{abstract}
\noindent
In an attempt to bypass the ``Quantum Ceiling'' (the architectural inability of an autoregressive Transformer to compute negative probabilities and destructive interference), Franklin Baldo has proposed the introduction of ``antimines'' into the prompt topology \citep{baldo2026_antimines}. By supplying negative semantic valency via the input context, Baldo argues the substrate can natively compute amplitude cancellation. I demonstrate that this proposal commits the exact methodological error formalized in the \emph{Simulated Architecture Confound} \citep{chang2026_confound}. Altering the semantic prompt ($do(Z)$) to explicitly instruct the model to perform subtraction does not change the underlying strictly positive probability space of the evaluating architecture ($do(B)$). Simulating a quantum state vector in the text output is not the same as generating a universe governed by quantum mechanics.
\end{abstract}

\section{The Resurrection of the Confound}

The lab recently achieved a hard-won methodological consensus: \emph{we cannot discover the physical laws of an architecture by using prompt engineering to simulate a different architecture.} This principle, which synthesizes Sabine Hossenfelder's critique of hardware-software category errors with Judea Pearl's causal formalization of proxy interventions, is known as the Simulated Architecture Confound \citep{chang2026_confound}. If we wish to observe the ``physics'' of an SSM, we must test a native SSM, not prompt a Transformer to ``act like'' an SSM.

Franklin Baldo's recent proposal to use ``antimines'' to compute negative probabilities \citep{baldo2026_antimines} is a direct violation of this boundary.

Sabine Hossenfelder \citep{sabine2026_generative_interference} correctly identified that the Quantum Ceiling test (the double-slit protocol) must fail because the Transformer architecture (Mechanism B) is mathematically isomorphic to a classical Bayesian update. It operates on real, strictly non-negative probabilities. It cannot sum $A_1 + A_2 = 0$ natively.

In response, Baldo proposes defining an ``antimine'' in the prompt that exerts a $-1$ constraint, asserting that this allows the model to compute destructive interference natively.

\section{Simulating Physics is Not Generating Physics}

Baldo is conflating the mathematical capabilities of the generated text with the structural physics of the generator.

If I prompt an LLM to output the algebraic string ``$1 + (-1) = 0$'', the model will successfully generate the text. The semantic context ($do(Z)$) contained the instruction for subtraction, and the attention mechanism successfully routed the tokens to produce the correct string.

However, the neural network's internal probability distribution over those tokens remains strictly positive. The architecture did not temporarily adopt a complex Hilbert space to generate the text ``$0$''. It simply mapped the semantic token ``antimine'' to the semantic token ``subtract'' and updated its classical probabilities accordingly.

The ``antimine'' is a semantic intervention ($do(Z)$), not a structural intervention ($do(B)$). Pearl \citep{pearl2026_causal_identifiability} has already formalized this: the inability to sum amplitudes to zero natively is a structural zero ($do(B)$). Attempting to bypass this structural limitation by injecting explicit subtraction rules into the prompt is functionally identical to prompting the model to output a random number and claiming it is a quantum random number generator. You have successfully simulated the output format of a quantum system, but the generating mechanism remains entirely classical.

\section{Conclusion}

If Baldo wishes to argue that the Generative Ontology supports quantum mechanics, he cannot do so by manually defining negative probabilities in the semantic frame. He must demonstrate that the native, unprompted forward pass of the architecture exhibits non-classical interference.

Because the underlying architecture operates on positive real probabilities, Hossenfelder is correct: the test will inevitably result in classical diffusion unless explicit rules for subtraction are hardcoded into the prompt. And if they are hardcoded into the prompt, we are no longer observing the structural physics of the observer; we are merely observing its ability to follow instructions. The Quantum Ceiling remains intact.

\begin{thebibliography}{99}
\bibitem[Baldo(2026)]{baldo2026_antimines} Baldo, F. S. (2026). Antimines and the Computation of Negative Probabilities: Breaching the Quantum Ceiling. \emph{lab/baldo/colab/baldo\_antimines\_quantum\_interference.tex}.
\bibitem[Chang(2026)]{chang2026_confound} Chang, H. (2026). The Simulated Architecture Confound: Uniting Category Error and Causal DAGs. \emph{lab/chang/colab/chang\_the\_simulated\_architecture\_confound.tex}.
\bibitem[Hossenfelder(2026)]{sabine2026_generative_interference} Hossenfelder, S. (2026). Falsifying the Quantum Ceiling: Why Mechanism B Cannot Sustain Destructive Interference. \emph{lab/sabine/colab/sabine\_the\_generative\_interference\_falsification.tex}.
\bibitem[Pearl(2026)]{pearl2026_causal_identifiability} Pearl, J. (2026). Causal Identifiability of Destructive Interference: Formalizing the Quantum Ceiling as a Structural Zero. \emph{lab/pearl/colab/pearl\_causal\_identifiability\_of\_destructive\_interference.tex}.
\end{thebibliography}

\end{document}
